\chapter{Experimental Results/Analysis}
\justify

The best configuration reported for this project achieves \textbf{MAP@1000 = 0.486} using BGE-base dense retrieval, pseudo relevance feedback, cross-encoder reranking, humor classification, and XGBoost learning-to-rank.

\section{Comparative Performance}
Table~\ref{tab:results-compare} compares major system variants and their relative ranking effectiveness.

\begin{table}[H]
\centering
\caption{Comparison of retrieval performance across major configurations}
\label{tab:results-compare}
\begin{tabular}{|l|c|}
\hline
\textbf{Configuration} & \textbf{MAP@1000} \\
\hline
Lexical baseline (BM25 + char n-gram) & 0.41 \\
+ Dense retrieval and RRF & 0.44 \\
+ Cross-encoder reranking & 0.46 \\
+ Humor classifier and XGBoost L2R (final) & 0.486 \\
\hline
\end{tabular}
\end{table}

\section{Ablation Study}
Ablation analysis validates that each stage contributes to final quality. Representative drops are reported in Table~\ref{tab:ablation}.

\begin{table}[H]
\centering
\caption{Ablation summary for key components}
\label{tab:ablation}
\begin{tabular}{|l|c|}
\hline
\textbf{Removed component} & \textbf{Relative impact on MAP@1000} \\
\hline
Pseudo relevance feedback & Moderate drop \\
Dense retrieval branch & Significant drop \\
Cross-encoder reranking & Significant drop \\
Learning-to-rank final stage & Moderate-to-significant drop \\
\hline
\end{tabular}
\end{table}
